% !TeX root = ../../Book3.tex
\section{零点}
\label{b3:sec:2605}

函数方程与 Euler 乘积共同限制零点，但边界直线上的无零点性需要额外论证。本节先精确区分已经排除的区域与下一章要处理的边界。

\subsection{平凡零点}
\label{b3:subsec:260501}

\begin{proposition}\label{b3:prop:26-trivial-zeros}
在 \(\operatorname{Re}s<0\) 中，\(\zeta\) 的零点恰为负偶整数，且都是单零点；此外 \(\zeta(0)=-1/2\)。
\end{proposition}
\begin{proof}
在该半平面，\(1-s\) 实部大于一，故 \(\zeta(1-s)\ne0\)，\(\Gamma(1-s)\) 全纯非零。函数方程中只有正弦因子可能为零，其零点为负偶整数，且均单。零处的值则直接由延拓积分的 \(-1/s\) 项乘以 \(s/2+O(s^2)\) 得到。
\end{proof}
这些零点称为\term[pingfan-lingdian]{平凡零点}[trivial zero]；“平凡”指其位置直接由显式因子确定，不表示相关极限可省略计算。

\subsection{临界带}
\label{b3:subsec:260502}

除负偶数外，其他零点只能位于闭条带 \(0\le\operatorname{Re}s\le1\)，因为右半平面一以右由 Euler 乘积排除。开条带 \(0<\operatorname{Re}s<1\) 称为\term[linjie-dai]{临界带}[critical strip]。\chapref{b3:ch:27}将证明右边界一上无零，再由函数方程排除左边界上非零虚部的零点。当前使用“闭条带”保留了这一尚需证明的步骤。

\subsection{临界线}
\label{b3:subsec:260503}

直线 \(\operatorname{Re}s=1/2\) 称为\term[linjie-xian]{临界线}[critical line]。在这条线上，消去极点后的对称函数将取实值；这使寻找中线零点可以化成实函数的变号问题。但是实值性不保证函数必变号，也不排除带内其他零点。

\subsection{非平凡零点}
\label{b3:subsec:260504}

不属于负偶数的零点称为\term[feipingfan-lingdian]{非平凡零点}[nontrivial zero]。它们的重数按全纯函数局部零点阶数定义；关于它们是否全部为单零点，本书不作假设。实区间 \((0,1)\) 上没有零点：交错级数
\(\eta(s)=\sum(-1)^{n-1}n^{-s}\) 在 \(\operatorname{Re}s>0\) 全纯，并满足 \(\eta=(1-2^{1-s})\zeta\)。对实 \(0<s<1\)，将相邻两项配对得到严格正数，分母 \(1-2^{1-s}<0\)，故 \(\zeta(s)<0\)。严格正性可由第一对正项及其余非负对看出。

\subsection{Riemann 猜想}
\label{b3:subsec:260505}

交错级数还说明一个表达式的分母零点未必是延拓函数的极点。在 \(\operatorname{Re}s>0\) 中，恒等式
\[
 \eta(s)=(1-2^{1-s})\zeta(s)
\]
先在绝对收敛半平面通过偶数项分离得到，再按恒等定理延拓。因子 \(1-2^{1-s}\) 在
\(s_k=1+2\pi\mathrm ik/\log2\) 处有单零点。若 \(k\ne0\)，\(\zeta\) 在该点全纯，所以 \(\eta(s_k)=0\)，商式中的表面奇点可以去除；只有 \(k=0\) 对应真正的极点，此时 \(\eta(1)=\log2\ne0\)。因此这个商可以用于跨过直线一，但必须处理这些离散的表面奇点，不能把分母的每个零点都报告为 zeta 的极点。

\ProseParagraphBreak

对于实 \(0<s<1\)，也可从相邻配对获得定量符号：第一对为 \(1-2^{-s}>0\)，其余各对为非负数，故
\[
 \eta(s)\ge1-2^{-s}>0,
 \qquad
 \zeta(s)\le-\frac{1-2^{-s}}{2^{1-s}-1}<0.
\]
这个不等式只在实区间使用了项的正性；复参数时各项没有这样的序关系。因此不能把实轴无零点论证直接平移成临界带内的无零点论证。

\ProseParagraphBreak

\begin{conjecture}[Riemann]\label{b3:conj:26-rh}
Riemann zeta 函数的每个非平凡零点都满足 \(\operatorname{Re}s=1/2\)。
\end{conjecture}
这称为\term[riemann-caixiang]{Riemann 猜想}[Riemann hypothesis]。它比素数定理所需的边界无零点性强得多。有限高度的数值验证只约束被检查的有限区域，即使精度和零点计数都严格，也不覆盖无穷高处的零点。本章随后证明零点整体上无限多，这同样不能替代它们全部位于中线的断言。
